\documentclass[conference]{IEEEtran}
\IEEEoverridecommandlockouts
\usepackage[T1]{fontenc}
\usepackage[utf8]{inputenc}
\usepackage{cite,amsmath,amssymb,graphicx,booktabs,array,tabularx}
\usepackage{algorithm,algorithmic}
\usepackage{tikz,pgfplots}
\usepackage[hyphens]{url}
\usepackage{xurl}
\usepackage[hidelinks]{hyperref}
\usepackage{microtype}
\usetikzlibrary{arrows.meta,positioning,calc}
\pgfplotsset{compat=1.17}
\definecolor{inkblue}{RGB}{42,97,130}
\definecolor{mutedgreen}{RGB}{71,123,103}
\definecolor{mutedamber}{RGB}{173,124,52}
\definecolor{mutedred}{RGB}{153,77,82}
\newcolumntype{Y}{>{\raggedright\arraybackslash}X}
\newcommand{\sys}{GDN Tree-Scan}
\newcommand{\pending}[1]{\textit{Planned; results pending: #1.}}
\begin{document}
\title{GDN Tree-Scan: Design and Optimization of Tree Verification for Recurrent-Hybrid Language Models}
\author{\IEEEauthorblockN{Zhiyuan Ma}
\IEEEauthorblockA{\texttt{zhiyuanma0520@gmail.com}}
\thanks{Working revision, September 24, 2026. Current production design and named coding-agent workloads; superseded implementations are archived.}}
\maketitle
\begin{abstract}
Tree speculative decoding for long-running coding agents must preserve coherent continuation across repeated generation and tool calls. We present GDN Tree-Scan, a Gated DeltaNet verifier that connects branch-local computation to accepted-path publication across recurrent, convolution, and attention state. Its implementation combines sequential scan and replay with fixed-shape device acceptance, fused draft selection, captured recurrent replay, and forked FlashAttention-2 with grouped split-K execution. On Qwen3.8-27B NVFP4 on GB10, among runs producing nonempty patches on both shared SWE-bench Verified Astropy tasks, the best recorded tree run achieves 29.09 pooled tokens/s versus SGLang EAGLE's 26.89, an 8.18\% higher rate. In a separate ten-task deployment, all tree-served agent attempts terminated and produced nonempty patches, with no degeneration detected by the recorded trace checks. Component gates check candidate-selection parity and deterministic bounded attention error for the deployed implementation.
\end{abstract}

\begin{IEEEkeywords}
speculative decoding, Gated DeltaNet, state commitment, numerical reproducibility, language-model serving
\end{IEEEkeywords}

\section{Introduction}
Speculative decoding reduces serial target-model invocations by verifying draft tokens in parallel. Its sampling guarantee assumes that verification supplies the intended target probabilities. Tree speculation expands the candidate set, but it does not relax that assumption~\cite{leviathan2023fast,chen2023sampling,miao2023specinfer}. For recurrent-hybrid models, realizing the verifier requires more than an attention mask: a candidate must inherit the state of its ancestors, and the next serving iteration must inherit exactly the path selected for continuation~\cite{yang2024gated,wu2025stree}.

\sys{} turns a candidate token tree into a target-model forward and a native continuation state. Its design connects an ancestry-aware attention path, branch-local GDN computation, a device acceptance walk, and accepted-path publication across recurrent, convolution, and attention key--value (KV) storage. A shared descriptor carries the logical tree through each stage~\cite{lumo2026archive,lumo2026stateless}. This system boundary matters because a verifier can compute plausible candidate logits yet publish the wrong state for the next iteration.

The central execution choice is to verify branches in temporary state and replay only the selected path into the native running state. The verifier scans branch paths, while the fixed-shape committer replays accepted operands through native GDN updates in a captured graph. Their finite-precision agreement is a separate requirement, not assumed from the recurrence. This choice exposes concrete optimization questions: how much branch state to cache, when to recompute ancestry, how to organize independent GPU work, and how to keep full-vocabulary probability arithmetic on the device. The attention layout and accepted-state mappings constrain those changes as much as the recurrent equation does.

\begin{figure}[t]
\centering
\input{figures/pipeline.tikz}
\caption{The tree verifier as a serving pipeline. A shared descriptor connects candidate generation, ancestry-aware target computation, device probability arithmetic, and accepted-path publication. The attention backend determines the physical mapping; all state surfaces must describe the same selected prefix.}
\label{fig:pipeline}
\end{figure}

The paper develops three implementation contributions. \textit{First}, it presents a fixed-shape tree verifier and continuation interface: candidates follow their ancestors during verification, and the selected materialized path becomes the next native state. \textit{Second}, it explains the GPU path-scan and replay design, including transient cut states, shared rounded operands, and publication into request-owned running rows. \textit{Third}, it describes the optimizations integrated into the deployed system: single-logits drafting, fused full-vocabulary candidate selection, spine-first attention layout, grouped-query split-K attention, and device acceptance and captured recurrent replay. The method follows the current fixed-32 Hydra27 implementation. Component gains require an ablation on the same agent workload.

The target workload is a long-running coding agent: generation alternates with repository inspection, tool execution, and tests, and a task can require many model requests. Performance therefore depends on useful completed work as well as token production. Our workload evaluation uses explicitly identified Astropy tasks from SWE-bench Verified~\cite{jimenez2024swe}. We report the model, serving route, task cohort, and measurement definition together; a decode-rate change is not a task-completion speedup. Component legality checks and loaded-kernel receipts support the implementation alongside the agent-workload measurements.

Tree verification and compact accepted-state reconstruction have direct prior art. Section~\ref{sec:related} positions our implementation against those mechanisms; we do not claim the first hybrid-tree system or a new GDN recurrence. The architecture and kernel design are developed in Sections~\ref{sec:contract}--\ref{sec:kernel}, and Section~\ref{sec:optimization} separates optimization mechanisms from their measured scope. Measurement provenance is documented in Appendix~\ref{app:audit}; superseded results and local-document timing numbers are excluded from the performance evaluation.

\section{Background and Related Work}
\label{sec:related}
\subsection{Drafting, verification, and serving}
Speculative sampling constructs an exact target sample from proposal and residual distributions, conditional on the probabilities actually supplied~\cite{leviathan2023fast,chen2023sampling}. SpecInfer supplies the tree-verification framing; Sequoia, EAGLE-2, and OPT-Tree investigate tree construction and its performance consequences~\cite{miao2023specinfer,chen2024sequoia,li2024eagle,wang2024opt}. Our study keeps these algorithmic questions separate from whether a deployed kernel and cache return the intended probabilities.

PagedAttention and SGLang address the scheduling, memory, and execution substrate~\cite{kwon2023efficient,zheng2024sglang}. Hybrid serving additionally manages recurrent-state pools alongside KV storage, as documented by SGLang's hybrid-model integration~\cite{pytorch2025sglang,sglang2026qwen}. Thus neither a fast isolated recurrent kernel nor low transient state memory alone determines user-visible latency. The deployed system uses forked FlashAttention-2 for tree attention, with native-compatible prefill and a specialized fixed-row decode kernel. Its physical layout and work partitioning are part of the verifier design~\cite{dao2023flashattention,lumo2026archive}.

\subsection{Recurrent and hybrid speculation}
Gated DeltaNet augments a decayed recurrent state with a data-dependent delta update~\cite{yang2024gated}. Parallel delta-rule formulations provide relevant algebraic background~\cite{yang2024delta}. Stateful speculation is established territory: STree studies speculative trees for hybrid state-space models, SpecMamba targets an FPGA implementation, and component-aware self-speculation studies how hybrid components affect the design~\cite{wu2025stree,zhong2025specmamba,borobia2026component}. These works motivate an explicit state lineage; they do not justify treating every recurrent update or numerical implementation as interchangeable.

\subsection{Direct GDN comparisons}
\textbf{Bole} derives a tree-structured closed form and evaluates it with a value-tiled finite-polynomial solver. It stores compact update factors and reconstructs the accepted state. Its SGLang integration combines batch-wide verification budgeting with captured GPU execution. Evaluation includes unquantized Qwen3.5-27B on GB10 and online replay of agent sessions~\cite{wang2026bole}. Its state and serving mechanisms directly overlap the problem studied here. The cost of our replay implementation is not a lower bound on alternatives.

Concretely, Bole's correction system is $(I+G)U=R$, where $G$ contains strict-ancestor interactions and $R$ is the source correction. Tree depth $d$ gives $G^{d+1}=0$, so $U=\sum_{m=0}^{d}(-G)^mR$. This finite identity changes the arithmetic schedule from sequential updates to matrix propagation~\cite{wang2026bole}. Our path preserves sequential updates and replays accepted operands; Section~\ref{sec:reconstruction-numerics} explains why the numerical comparison remains an empirical question.

\textbf{TreeWY} expresses GDN tree verification as a triangular system and reconstructs the accepted state from pseudo-values. Its vLLM evaluation uses Qwen3.5-35B and 397B on B200, greedy verification, and disabled prefix caching. It explicitly permits finite-precision disagreement with the baseline and reports that wider trees improve acceptance without yet winning throughput~\cite{ghantasala2026treewy}. Its numerical contract and hardware differ from the deployed system studied here.

\textbf{ReplaySSM} caches recurrent inputs and changes when state is materialized; its implementation RFC includes GDN and Mamba2~\cite{replayssm2026}. Reconstruction and replay are not claimed as new concepts here. Table~\ref{tab:related} states the comparison boundary. Published speed ratios from different checkpoints, precision formats, workloads, and software builds are not combined into a performance ranking.

\begin{table}[t]
\caption{Direct mechanism-comparison scope. The authors' implementations of the three external methods listed here were not executed in this study.}
\label{tab:related}
\centering\footnotesize
\begin{tabularx}{\columnwidth}{@{}>{\raggedright\arraybackslash}p{0.22\columnwidth}YY@{}}
\toprule
Method & Overlapping mechanism & Boundary for this study \\
\midrule
Bole~\cite{wang2026bole} & Parallel verifier; factorized commit & Different stack and precision; no matched run \\
TreeWY~\cite{ghantasala2026treewy} & Tree solve; accepted-state reconstruction & Different models and operating regime; no matched run \\
ReplaySSM~\cite{replayssm2026} & Input caching; deferred state work & Alternative state policy; no matched run \\
\sys{} & Scan/replay; accepted-path remapping & Deployed fixed-shape verifier; route-bound checks and agent workloads \\
\bottomrule
\end{tabularx}
\end{table}

\section{Tree-Verifier Architecture}
\label{sec:contract}
\subsection{State and the execution boundary}
Let $m$ be a materialized token prefix: the tokens already consumed by the target's forward computation. Write its persistent state as
\begin{equation}
\mathcal{S}(m)=\bigl(R(m),C(m),K(m)\bigr),
\end{equation}
where $R$ collects recurrent matrices, $C$ collects convolution histories, and $K$ collects attention KV entries. The logical emitted prefix can also include a newly sampled token $b$ that has not yet been consumed. The continuation boundary is then $(\mathcal{S}(m),b)$, not a claim that $\mathcal{S}(mb)$ has already been computed. This distinction is essential when counting the correction or bonus token while checking cache state.

A tree descriptor contains candidate token IDs, parent indices, depth/position information, attention visibility, and maps between logical nodes and physical slots. For a candidate node $i$, each target layer must consume the state obtained from the candidate's root-to-parent path. At publication, the selected materialized path, pending-token convention, and next-drafter input row must agree. Figure~\ref{fig:contract} illustrates the three state surfaces documented by the implementation audit~\cite{lumo2026stateless,lumo2026archive}.

\subsection{The recurrence and path locality}
For one GDN head, let $R\in\mathbb{R}^{d_v\times d_k}$. Queries and keys are in $\mathbb{R}^{d_k}$ and values in $\mathbb{R}^{d_v}$. With decay $\alpha_i$ and write strength $\beta_i$, a node applies the standard update~\cite{yang2024gated}
\begin{align}
R'_i &= \alpha_i R_{\pi(i)}, &
\delta_i &= \beta_i\bigl(v_i-R'_i k_i\bigr),\\
R_i &= R'_i+\delta_i k_i^\top, &
o_i &= R_iq_i,
\end{align}
where $\pi(i)$ is its parent. A previous node in the packed buffer need not be $\pi(i)$. The causal convolution likewise requires path predecessors rather than adjacent packed rows. The attention visibility relation must select the same logical prefix and ancestry. These requirements follow from the model's computation; a correct attention mask cannot repair a wrong recurrent parent.

The implementation verifies nodes with branch-local scratch and makes the selected continuation durable through replay and remapping. This is an implementation route, not the only way to satisfy the contract~\cite{lumo2026archive,lumo2026stateless}. In particular, rejecting a branch requires that its contents cannot influence future computation; it does not require zeroing every scratch byte if future reads are correctly restricted.

\subsection{A descriptor shared by verification and commitment}
Figure~\ref{fig:pipeline} shows the served path. Native multi-token prediction (MTP) supplies the principal draft chain; selected variants add runner-up candidates. The descriptor maps each candidate through the attention mask, recurrent scan, sampler, and selected-path publication. The agent deployment uses Hydra27: 27 active draft nodes excluding the root, packed into 31 physical draft slots plus the root (32 verification rows). Its five-depth MTP head contributes 15 nodes; suffix continuation contributes a six-token principal tail and four- and two-token rescue branches. The maximum draft depth is eleven; four post-root MTP forwards construct the head. Four masked physical slots are padding, not extra logical candidates. Suffix candidates come from the configured Arctic suffix predictor. The current implementation derives the mask, physical packing, acceptance walk, and GDN path schedule from this topology~\cite{lumo2026archive}.

\begin{figure}[t]
\centering
\input{figures/state-contract.tikz}
\caption{The accepted-prefix interface. Recurrent, convolution, and attention-KV state must all describe the same materialized prefix. A newly sampled correction or bonus token can remain pending until the next forward.}
\label{fig:contract}
\end{figure}

The FA2 route places the principal spine first in the physical suffix and permutes KV write locations and ancestry-mask key columns while retaining logical query-row and position order. Selection remains in logical-node space; accepted-path publication maps those nodes back to their physical slots. This distinction prevents a layout optimization from silently changing which key/value entries a branch reads. Recurrent paths begin at the native pre-tree state or a transient cut state; the device committer publishes the selected recurrent and convolution state into native running rows and re-linearizes its attention KV entries~\cite{lumo2026archive,lumo2026stateless}.

\begin{algorithm}[t]
\caption{Logical verifier and continuation boundary}
\label{alg:verify}
\begin{algorithmic}[1]
\REQUIRE Reference-aligned state boundary, candidate descriptor $T$
\STATE Map logical candidates to physical slots and positions.
\STATE Verify using path-local convolution, GDN, and attention.
\STATE Apply the configured sampler to matching target/proposal rows.
\STATE Obtain the accepted path and any pending correction token.
\STATE Publish recurrent and convolution state for the materialized path.
\STATE Map its attention KV entries to the native continuation layout.
\STATE Map the selected hidden-state row to the next drafter input.
\STATE Preserve the native pending-token and position convention.
\RETURN Continuation boundary and emitted-token accounting
\end{algorithmic}
\end{algorithm}

\section{GPU Scan and Accepted-Path Replay}
\label{sec:kernel}
\subsection{Path-parallel verification and state lifetime}
The deployed GDN verifier decomposes the tree into dependent path groups. A first path starts from the request's pre-tree state and exports the cut states needed by downstream paths. The second level evaluates independent branch paths from those cut states. Each path applies sequential node updates in ancestry order while separate programs cover value tiles and heads. Candidate outputs continue through the model; rounded operands are recorded for later accepted-path replay~\cite{lumo2026archive}.

Cut states and operand rings are temporary within a verification step. The durable continuation resides in the request-owned native recurrent and convolution running rows. Replay starts from the pre-tree running state and publishes the accepted frontier back to those rows; prefix-cache snapshots and restores use that native layout. Static work buffers can be reused by CUDA graphs, but they do not replace the request's persistent state with a branch-leaf pointer. The current lifecycle does not depend on explicitly zeroing every scratch buffer after each step.

This distinction also separates the deployed path-group implementation from an experimental single-launch GDN kernel. The latter removes the intermediate cut-state exchange but is disabled in the reported agent deployment. We therefore neither describe the deployed verifier as one GDN launch nor claim that it eliminates all temporary state traffic.

\begin{algorithm}[t]
\caption{Logical recurrence and selected-path publication. The implementation groups independent paths into GPU launches and retains only their needed cut states.}
\label{alg:scan-replay}
\begin{algorithmic}[1]
\REQUIRE Pre-tree state $R_0$, descriptor $T$, rounded node operands $x_i$
\FOR{node $i$ in topological order}
\STATE $R_p\gets R_0$ if $i$ has no candidate parent; otherwise $R_p\gets H[\pi(i)]$
\STATE $(H[i],o_i)\gets\operatorname{VerifyStep}(R_p,x_i)$
\ENDFOR
\STATE Finish target forward and select accepted materialized path $A$
\STATE $R\gets R_0$
\FOR{node $i$ in $A$, in path order}
\STATE $(R,\_)\gets\operatorname{NativeUpdate}(R,x_i)$
\ENDFOR
\STATE Publish $R$ and matching convolution history to native running rows
\STATE Remap accepted KV entries and retire speculative scratch
\RETURN Matching continuation state and any pending token
\end{algorithmic}
\end{algorithm}

\subsection{Sequential arithmetic and accepted-path replay}
Verification uses sequential path scans with fp32 carried state, raw gate inputs, and in-kernel query/key normalization. The deployed fixed-32 committer uses native fused GDN updates with recorded accepted operands and in-place publication to the fp32 state bank. Its graph contains one native update call per GDN layer. Query/key normalization, beta conversion, gate evaluation, and output casts can differ between these implementations; agreement is not guaranteed by their common real-arithmetic recurrence~\cite{lumo2026archive}.

The accepted path is a device product of the sampling stage. A captured recurrent committer consumes its node indices and lengths and publishes recurrent state. Separate convolution and attention-KV publication consume the same selected path. One graph replay launches the sequence of forty-eight per-layer native updates. It is not a single fused all-layer kernel and does not establish bit identity between verification and publication arithmetic.

\subsection{Temporary state versus recomputation}
Path decomposition caches the states needed where downstream paths branch, then reuses them across independent work. If a program exports $N_c$ cut states with value-tile width $B_v$ and key dimension $d_k$, their fp32 payload is
\begin{equation}
 M_{\mathrm{cut}}=4N_cB_vd_k\ \text{bytes}.
\end{equation}
This is a logical payload, not a register-allocation or total-memory measurement. Recomputing ancestry could remove some exports at the cost of repeating updates. Exporting every candidate state would instead trade more storage for less recomputation. The current two-level path schedule occupies a particular point in this design space.

Accepted-path replay carries one state tile through a bounded path loop. Its carried-state footprint is $O(B_vd_k)$ in tree size, while the recorded operands and candidate outputs scale with the verification shape. Static padding makes graph buffers and kernel geometry reusable, but padded rows must remain invisible to attention and selection. Neither a shorter logical path nor fewer state writes alone establishes a faster complete request.

\begin{figure}[t]
\centering
\input{figures/scan-replay.tikz}
\caption{State lifetime in the deployed path-scan verifier. Temporary cut states connect path groups; recorded operands support selected-path replay. Only the accepted frontier becomes native persistent state.}
\label{fig:scan-replay}
\end{figure}

Compact triangular-solve and finite-Neumann methods choose another point in this design space: they form correction factors and reconstruct state without replaying every accepted rank-one update. Their arithmetic schedule differs from sequential scan and replay. Lower state storage or less replay work alone does not establish lower complete-step latency, and a real-arithmetic identity does not settle finite-precision agreement~\cite{ghantasala2026treewy,wang2026bole}.

\section{Optimization Mechanisms and Backend Constraints}
\label{sec:optimization}
\subsection{Reuse logits and fuse full-vocabulary selection}
The deployed drafter uses the full vocabulary. At each head depth, a single logits tensor supplies the principal token and runner-up candidates. The principal chain advances through four captured MTP iterations after the root; sibling proposals reuse the head scores without independent drafter forwards. The suffix predictor extends selected paths within the same fixed verification capacity~\cite{lumo2026archive}.

A fused CUDA selection kernel emits both the spine token and the three ordered candidate indices directly into static graph buffers. It replaces separate argmax/top-k reductions and their output copies; it does not remove the vocabulary projection itself. Matching the pinned reference requires preserving two tie conventions: argmax chooses the lowest token index among equal maxima, while top-k has its own ordering among selected ties. The fused kernel reproduces both conventions, including duplicate proposals when the reference produces them. Changing that policy would change the proposal tree, so it is a separate algorithmic decision.

\subsection{Keep acceptance and publication on device}
The fixed-32 tree-acceptance walk consumes full-vocabulary target/proposal probabilities and produces accepted node indices, emitted tokens, path lengths, and the pending correction or bonus token. It runs through the fixed-shape device route, not the older generic Python node walk. Device-side product filling feeds the captured committer, avoiding host construction of per-layer path lists between selection and publication~\cite{lumo2026archive}.

Recurrent replay, convolution publication, attention-KV remapping, and the next drafter's hidden-row selection all consume that same path. CUDA graphs reduce repeated launch orchestration; they do not alter the required selected-prefix contract. Nor does a captured committer imply that every kernel within it has been fused into a single launch.

\subsection{Prepare convolution and replay operands for the fixed shape}
Tree convolution gathers a causal window for each candidate from its ancestors and the prior native convolution state. The deployed route uses static path-index preparation, fused tree convolution, and stacked operand rings spanning the GDN layers. Ordered tap operations and cast boundaries remain part of its arithmetic contract. The fixed-32 committer selects the accepted convolution window and writes it to the native running row; graph scratch is prepared for the next iteration rather than becoming durable per-node history~\cite{lumo2026archive}.

\subsection{Spine-first layout and grouped split-K attention}
\label{sec:reconstruction-numerics}
The forked FA2 kernel applies tree visibility inside attention. Its spine-first packing keeps the principal path contiguous in the candidate suffix. This matters numerically as well as logically: inserting masked branch columns between spine keys can change the reduction layout of an online softmax. Permuting KV slots and the matching ancestry-mask key columns preserves logical query order while controlling the physical key layout~\cite{lumo2026archive}.

The deployed GQA-pair kernel shares staged key/value tiles across two query heads. Split-K divides the context walk into four partitions. Each partition computes a local maximum, denominator, and weighted-value sum; a combine kernel rescales and merges these partial results. This exposes more independent work at batch one, but changes floating-point accumulation order. Its qualification therefore uses explicit determinism and numerical bounds against both the incumbent and a high-precision reference, rather than asserting byte identity with the incumbent.

The workload receipt identifies the specialized kernel in all sixteen full-attention layers and records the binary hash, four-way split, graph identity, and no-fallback dispatch. Prefill uses the patched FA2 integration's native-compatible route. These execution records establish which implementation generated the workload observations; performance attribution still requires a matched component ablation.

\subsection{Publish one coherent native continuation}
If accepted logical node $a_j$ becomes continuation position $j$, the KV remapper copies from the physical slot assigned to $a_j$ into the request's native destination slot for $j$. The slot permutation is applied to verification reads and writes as well as publication. A valid ancestry mask alone cannot repair inconsistent addressing.

Recurrent and convolution publication make the same selected frontier visible in native running rows. Prefix caching and full/piecewise CUDA graphs are enabled in the reported deployment. Their state ownership follows the native request lifecycle; temporary cut-state exports remain a distinct within-step resource.

\begin{table*}[t]
\caption{Optimizations in the recorded current deployment. Work removed is a mechanism description; no component timing is converted into an agent-workload speedup.}
\label{tab:optimization-map}
\centering\footnotesize
\begin{tabularx}{\textwidth}{@{}p{0.20\textwidth}YYY@{}}
\toprule
Mechanism & Work addressed & Deployed realization & Evidence boundary \\
\midrule
Single-logits drafting & Repeated projection and score preparation & Full vocabulary; four captured MTP iterations after root & Source and boot engagement \\
Fused top-3 selection & Separate argmax/top-k reductions and copies & One CUDA selection into graph buffers & Byte and graph-replay component gates \\
Path-parallel GDN scan & Serial branch work and repeated ancestry & Two-level path schedule with transient cut states & Source-selected route; single-launch candidate disabled \\
Device acceptance and commit & Host path construction and launch orchestration & Fixed-32 walk, device products, captured committer & Boot and route receipts; no blanket sampler-law proof \\
Convolution and operand preparation & Repeated path-index and per-layer preparation & Fused convolution, pregather, stacked layer rings & Enabled/engaged deployment path \\
Spine layout and KV remap & Inconsistent physical ancestry and continuation & Spine-first permutation plus accepted-path remap & Common slot mapping in verification/publication \\
FA2 grouping and split-K & Repeated KV staging and insufficient independent work & Query-head pairs; four context partitions and combine & Numerical credential matches the loaded binary \\
Native-state lifecycle & Persistent branch-specific cache ownership & Publish native running rows; graphs and prefix cache enabled & Current state contract and recorded agent behavior \\
\bottomrule
\end{tabularx}
\end{table*}

\section{Evaluation Setup}
\label{sec:protocol}
Performance is evaluated on coding-agent workloads, with the task subset and serving configuration attached to every reported rate. The Qwen3.8-27B NVFP4 evidence uses SWE-bench Verified Astropy tasks on a single GB10/DGX Spark: a shared two-task comparison of three tree runs against SGLang EAGLE, plus a separate ten-task tree deployment~\cite{lumo2026archive,jimenez2024swe}. The best recorded shared-task rate supplies the headline; later runs and the larger deployment remain in the evaluation. These selected subsets are not the full benchmark. Every decode-rate result uses the pooled completed-request estimator in Eq.~\ref{eq:pooled-decode}, excluding first-output latency. Recorded agent duration includes the agent loop, model service, tools, and repository operations; server startup and subsequent evaluator time are separate.

The system description follows the latest deployed Hydra27 route supported by these records. Component validation is tied to the fused-selection and split-K artifacts used by that route. Archived experiments on superseded implementations do not supply results or qualification claims for the current system. Newer topology or kernel candidates without a completed serving record are likewise excluded from the deployed-method description.

\section{Results}
\label{sec:results}
\subsection{SWE-bench Verified coding-agent workload}
\label{sec:agent-workload}
\input{results/agent-workload/case-study.tex}

\subsection{Validation of the deployed components}
\label{sec:current-validation}
The fused full-vocabulary selection gate compares the two output tensors against the pinned argmax and top-k reference operations. Across 1,368 input cases and five block-count settings (6,840 configurations), it records zero byte mismatches; planted negative controls fire. A separate captured four-level selection graph completes 24 replays with zero mismatches. These results validate candidate-selection outputs for the tested inputs and graph behavior, not end-to-end task equivalence~\cite{lumo2026archive}.

The split-K numerical credential is bound to the same FA2 binary and patch identity recorded by the ten-task deployment. It reports bit-identical repeated output and log-sum-exp values within each of sixteen determinism cases, with matching digests across two processes. On synthetic query/key/value tensors drawn at scales measured from banked model operands, 93.31\% of output elements agree with the incumbent within two ULP, the maximum absolute output difference is 0.00390625, and maximum log-sum-exp disagreement is four ULP. All nine declared checks pass, including comparisons against a float64 dense reference and checks for non-finite disagreements. These are attention-component bounds; split-K changes rounding and is not claimed to be byte-identical to the incumbent or to preserve every downstream token.

Loaded engagement is checked separately from component eligibility. The deployment receipt records the specialized FA2 arm in all sixteen attention layers, four splits, and no fallback, while the boot records identify fused selection, fixed-32 device acceptance, and captured state publication. The agent traces and evaluator outcomes in Section~\ref{sec:agent-workload} supply the complementary application evidence. Neither passing component gates nor observing no degeneration in a finite task cohort establishes a full-model conditional-distribution proof.

\section{Discussion and Remaining Evidence}
\label{sec:planned}
The deployed design couples tree geometry, physical layout, and continuation publication. Optimizing any one of these in isolation can move work elsewhere or invalidate another stage's mapping. The current implementation uses transient cut states, static graph buffers, and native persistent rows for different purposes; conflating them would obscure both memory cost and state ownership.

Sequential replay and compact reconstruction remain alternative implementation choices. The present paper does not benchmark the authors' Bole or TreeWY systems, and its deployed-component gates do not rank those methods. A direct comparison requires compatible models, matched workloads, and each system's actual execution path.

For long-running agents, task completion and task wall time are the relevant application outcomes. More decoded tokens can reflect additional reasoning or a longer tool loop. The pooled decode comparison is therefore accompanied by task identities, outcomes, agent/tool policy, and terminal budgets. It describes the evaluated deployments and the stated retrospective patch-producing subset; it is not a task-completion speedup estimate.

Current-component legality and whole-task performance answer different questions. The selection gate checks exact candidate outputs; the split-K gate checks deterministic bounded rounding; the workload tests exercise the composed system. A broader quality claim requires held-out tasks and repetitions. Isolating an optimization's application benefit requires toggling that component within the same eligible serving configuration and task protocol.

\section{Conclusion}
\sys{} implements tree verification for a recurrent-hybrid model by connecting an ancestry-aware target forward to accepted-path state publication. A shared descriptor coordinates attention visibility, branch-local GDN computation, convolution history, device probability arithmetic, and the next native continuation. Sequential scan and accepted-path replay make the state lifetime and arithmetic choices explicit; cached state, recomputation, physical layout, and independent-work batching expose the main optimization tradeoffs.

On two shared SWE-bench Verified Astropy tasks with Qwen3.8-27B NVFP4 on GB10, among runs producing nonempty patches on both tasks, the best recorded tree run attains 29.09 pooled tokens/s versus SGLang EAGLE's 26.89, an 8.18\% higher rate. Later tree runs on the same pair attain 28.28 and 27.27; the separate ten-task deployment attains 25.63. All ten attempts in that larger deployment terminated with nonempty patches and no degeneration detected by the recorded trace checks. The common rate definition compares the complete decoding systems; draft budgets, selection scope, and task outcomes are reported alongside the rates.

The contribution is a concrete verifier design and its implementation mechanisms, supported by current-component checks and workload-specific observations. A broader performance claim requires additional tasks and repeated runs with a declared tuning protocol. The present results do not establish broader task-quality preservation, full-model sequential equivalence, or general superiority over SGLang or the published compact-state systems.

\appendices
\section{Reproducibility and Claim Boundaries}
\label{app:audit}
The artifact records source revisions, effective serving settings, component binary identities, topology, graph/cache configuration, and original task metadata. The latest deployed design is determined by promotion history and positive execution receipts, not by the date of an experiment that happened to revisit an older route. Superseded implementations and their diagnostic results remain available as dated audit records; they are not presented as the current method.

All workload rates use the completed-request pooled decode estimator in Eq.~\ref{eq:pooled-decode}. Token totals and latency sums come from the same completed-request population. First-output latency, tool execution, and evaluator time are separate quantities. Earlier global-token/subset-wall proxies and local-document timing do not supply paper performance claims. The selection ledger preserves every included and excluded run-pair and the unchanged unfiltered measurements.

\textit{Algebraic correctness} concerns the real-arithmetic recurrence. \textit{Numerical agreement} concerns particular kernels and references. \textit{Sampler correctness} concerns the acceptance/residual law for the supplied probabilities. \textit{Serving correctness} additionally requires coherent state across continuation and cache operations~\cite{yang2024delta,chen2023sampling,lumo2026archive}. Equal candidate outputs, bounded attention error, and observed task completion address different parts of that chain.

For target probabilities $p$ and proposal probabilities $q$, ordinary rejection sampling combines accepted mass $\min(p,q)$ with residual mass proportional to $(p-q)_+$. This gives the intended target marginal only for the probabilities actually supplied~\cite{leviathan2023fast,chen2023sampling}. A correct abstract sampling rule cannot compensate for stale state or different target logits. The paper therefore does not promote component checks or task outcomes into a general distribution-preservation claim.

\section{Remaining Workload Comparisons}
\label{app:validation}
A matched current native-MTP/tree comparison would freeze the checkpoint, task cohort, agent harness, sampling, tool policy, and budgets, then report outcomes and completion time over every attempted task. Each method may use its own speculative depth and serving optimizations; the tuning protocol and selected configuration must be explicit. Confirmation should predeclare performance eligibility and retain unsuccessful attempts in the outcome report.

Further component ablations, concurrent-agent experiments, and Bole/TreeWY system comparisons are conditional additions to that protocol. They are not inferred from archived snippet timing, unserved kernel candidates, or experiments on superseded routes.

\label{ReferencesStart}
\bibliographystyle{IEEEtran}
\bibliography{ref}
\end{document}
